\section{Conclusiones}

\subsection{Resumen de la importancia del problema y de la necesidad de su estudio}
La presente investigación subraya la necesidad imperativa de trascender las simplificaciones lineales en el modelado de redes logísticas interregionales. La importancia del problema radica en que el comportamiento real de los costos de transporte y almacenamiento está intrínsecamente ligado a fenómenos no lineales, como las economías de escala y la saturación de infraestructura. Ignorar estas dinámicas no solo conduce a soluciones subóptimas, sino que genera una desconexión entre la teoría matemática y la realidad operativa. El estudio de este problema es esencial para proporcionar un marco analítico robusto que permita capturar la complejidad de los sistemas logísticos modernos, garantizando que la optimización se base en representaciones matemáticas fieles a la naturaleza física y económica del flujo de bienes.

\subsection{Conexión entre el planteamiento del problema y los objetivos de la investigación}
Existe una correspondencia directa y necesaria entre el planteamiento del problema y los objetivos definidos. Ante la brecha identificada entre los modelos lineales tradicionales y la realidad no lineal, la investigación se ha planteado objetivos que abordan esta problemática desde tres frentes complementarios:
\begin{enumerate}
    \item La construcción de un modelo algebraico riguroso responde directamente a la necesidad de formalizar las dinámicas de costo no lineales sobre una estructura de grafo.
    \item El análisis topológico y la derivación de las condiciones de Karush-Kuhn-Tucker (KKT) permiten fundamentar teóricamente la viabilidad de encontrar soluciones óptimas, proporcionando el rigor matemático que el planteamiento del problema exige.
    \item Finalmente, la implementación y evaluación de algoritmos numéricos de optimización continua cierran el ciclo de investigación, transformando el problema abstracto en una herramienta de resolución capaz de validar la eficiencia y convergencia del modelo propuesto. 
\end{enumerate}
De esta manera, el cumplimiento de los objetivos garantiza una respuesta integral a las preguntas de investigación planteadas al inicio del estudio.
